\documentclass[11pt]{article}

\usepackage[margin=1.2in]{geometry}
\usepackage{amsmath,amssymb,amsthm,mathtools}
\usepackage{hyperref}
\usepackage{setspace}
\usepackage{parskip}
\usepackage{titlesec}
\usepackage{bm}
\usepackage{graphicx}

\setstretch{1.15}
\setlength{\parindent}{0pt}

\newtheorem{theorem}{Theorem}[section]
\newtheorem{lemma}[theorem]{Lemma}
\newtheorem{proposition}[theorem]{Proposition}
\newtheorem{corollary}[theorem]{Corollary}
\newtheorem{conjecture}[theorem]{Conjecture}
\theoremstyle{definition}
\newtheorem{definition}[theorem]{Definition}
\newtheorem{remark}[theorem]{Remark}

\DeclareMathOperator{\tr}{tr}
\DeclareMathOperator{\diag}{diag}

\title{Constraint Closure and Entropic Decimation:\\
Toward a Unified Field Theory of Computation, Cognition, and Cosmology}
\author{Flyxion}
\date{\today}

\begin{document}

\maketitle

\begin{abstract}
Contemporary models of computation and cognition are dominated by predictive
paradigms that optimize local consistency without guaranteeing global
realizability. This essay develops an alternative framework in which physical
fields, cognitive processes, and computational systems are understood as
instances of constraint accumulation followed by entropy-respecting decimation.
Building on the Integral Decimation Method and tensor-train representations, we
reinterpret inference as a process of reconstructing globally admissible
configurations through sequential application of local operators and nonlinear
compression.

Within this framework, the Relativistic Scalar--Vector Plenum (RSVP) theory is
formalized as a field system in which a scalar field encodes constraint density,
a vector field encodes directional organization, and an entropy field regulates
admissibility rather than participating symmetrically in dynamics. The resulting
action functional introduces an explicit entropy floor dependent on local field
intensity, producing a soft but enforceable admissibility condition.

A computational scheme is derived by discretizing the action and applying a
Trotter-split decomposition into structural and admissibility gates, yielding a
tensor-train evolution governed by sequential local updates and singular value
truncation. This decimation process is interpreted as an entropy-filtering
mechanism that preserves globally consistent structure while eliminating
incoherent modes.

We propose a convergence conjecture stating that the resulting truncated sweep
map is contractive under appropriate conditions, supported empirically by
stabilization of bond dimension across sweeps. Categorical, information-theoretic,
dynamical-systems, and geometric analyses of the scheme are developed, connecting
it to sheaf-theoretic gluing, Landauer's principle, projected gradient flow, and
manifold curvature. This establishes a unifying principle: stable structure in
physics, cognition, and computation arises not from prediction but from constraint
closure under admissibility filtering. The essay thereby connects tensor-network
numerics, thermodynamic field theory, and models of intelligence within a single
geometric and informational framework.
\end{abstract}

\newpage
\tableofcontents
\newpage

% ============================================================
\section{Introduction: From Prediction to Reconstruction}
% ============================================================

The dominant paradigm in contemporary artificial intelligence treats
intelligence as prediction. Systems are trained to estimate the next token, the
next pixel, or the next state in a sequence, and performance is measured by the
accuracy of these local estimates. This approach has produced remarkable
practical successes, yet it exhibits a fundamental structural limitation.
Prediction ensures local plausibility, but it does not guarantee global
consistency. A sequence of locally plausible steps may fail to assemble into a
coherent whole.

This limitation is not merely technical but conceptual. Predictive systems
operate on fragments of structure without enforcing the constraints that make
those fragments mutually compatible. The resulting outputs often exhibit
internal contradictions, discontinuities, or failures of realizability that
cannot be repaired by improving local accuracy alone. The problem is not
insufficient prediction, but the absence of a mechanism that enforces global
admissibility.

In contrast, physical systems do not operate by prediction. They evolve under
constraints. A configuration of a physical field is not accepted because it is
locally plausible, but because it satisfies a set of governing relations
simultaneously. The global state of the system is therefore a solution to a
constraint problem rather than a sequence of guesses. Stability emerges when
these constraints are satisfied in a self-consistent manner.

This essay develops a framework in which computation, cognition, and physical
dynamics are unified under this constraint-first perspective. The central claim
is that intelligence and structure arise from the process of reconstructing
globally admissible configurations through the accumulation of local constraints
and the elimination of incoherent degrees of freedom. This process is formalized
as a combination of local operator application and entropy-respecting decimation.

Recent advances in tensor-network methods, particularly the Integral Decimation
Method of Grimm, Staat, and Eaves~\cite{grimm2026}, provide a concrete
computational realization of this idea. High-dimensional integrals, which are
intractable when treated as monolithic objects, can be decomposed into sequences
of local transformations followed by controlled compression. The key insight is
that the exponential weight $e^{iS[\psi]}$ over an action $S$ can be
reconstructed as a product of gate operators acting sequentially on an initially
unentangled state, changing the computational complexity of integration from
exponential to polynomial~\cite{grimm2026, oseledets2011tt}.

At the same time, the RSVP framework provides a field-theoretic language in
which constraint density, directional organization, and entropy are represented
explicitly. By combining these perspectives, we obtain a unified description in
which the evolution of a system is governed by local interactions while global
consistency is enforced through decimation.

The key shift is therefore from prediction to reconstruction. Rather than asking
what the next state should be, we ask which configurations can exist at all
under the imposed constraints. The computational task becomes one of assembling
a globally consistent structure from local components while discarding those
components that cannot be integrated into a coherent whole. This shift aligns
naturally with tensor-network representations, in which a global object is
never stored explicitly but assembled from locally structured cores whose
internal dimensions --- the bond dimensions --- measure the retained
correlations between neighboring regions.

This perspective has implications beyond numerical methods. It suggests that
cognition itself may operate as a process of constraint closure, in which
hypotheses are generated and then filtered by their compatibility with an
evolving internal model. It also suggests that physical structure arises not
from expansion or accumulation, but from the progressive elimination of
inadmissible configurations.

The remainder of the essay develops this framework in detail. Section~2
formalizes the RSVP action with entropy as a regulator of admissibility,
deriving the equations of motion and establishing the asymmetric role of the
entropy field. Section~3 derives a prototype decimation scheme based on a
Trotter-split decomposition of the action, leading to a tensor-train
implementation with exact commutativity of the Trotter factors on the
local basis. Section~4 describes the implementation and the observables
extracted from it, focusing on bond-dimension diagnostics as signatures
of constraint closure. Section~5 presents empirical evidence for convergence
and formulates the convergence conjecture, relating it to the renormalization
group. Sections~6--9 develop categorical, information-theoretic, dynamical-systems, and geometric perspectives on constraint closure. Section~10 surveys extensions and open problems. Section~11 concludes.

% ============================================================
\section{The RSVP Action with Entropy as Regulator}
% ============================================================

\subsection{Field Variables and Their Asymmetric Roles}

Let $\Omega \subset \mathbb{R}^d$ be a compact domain. We define the RSVP
state as a triple
\begin{equation}
  X = (\Phi, \mathbf{v}, S) \in C^\infty(\Omega) \times
      C^\infty(\Omega;\mathbb{R}^d) \times
      C^\infty(\Omega;\mathbb{R}_{\ge 0}).
\end{equation}

We distinguish the three RSVP fields by function rather than by symmetry. The
scalar field $\Phi : \Omega \to \mathbb{R}$ encodes local constraint density or
potential burden. It measures how strongly a region of the domain is required to
satisfy structural relations. The vector field $\mathbf{v} : \Omega \to
\mathbb{R}^d$ encodes directional transport or organizing flow. It carries
information about how structure propagates, aligns, or redistributes across the
domain.

The entropy field $S : \Omega \to \mathbb{R}_{\ge 0}$ plays a fundamentally
different role. It does not participate symmetrically in the dynamics alongside
$\Phi$ and $\mathbf{v}$. Instead, it regulates the admissibility of
configurations produced by them. A configuration of $(\Phi,\mathbf{v})$ is not
automatically allowed; it must be supported by sufficient entropy to remain
stable and realizable. In this sense, $S$ functions as a certification field,
determining which configurations can exist rather than contributing directly to
their formation.

This asymmetry is essential. If $S$ were treated as an ordinary dynamical field
on equal footing with $\Phi$ and $\mathbf{v}$, the distinction between
structure and admissibility would collapse. The purpose of the RSVP formulation
is precisely to preserve this distinction.

\begin{definition}[Admissibility]
  A configuration $X$ is \emph{admissible at $x \in \Omega$} if
  \begin{equation}
    S(x) \ge S_*(\Phi(x), \mathbf{v}(x)),
    \label{eq:admissibility}
  \end{equation}
  where $S_* : \mathbb{R} \times \mathbb{R}^d \to \mathbb{R}_{\ge 0}$ is the
  \emph{entropy floor function}. The configuration is globally admissible if
  \eqref{eq:admissibility} holds for almost every $x \in \Omega$.
\end{definition}

\subsection{The Entropy Floor}

We define the entropy floor as a quadratic in local field intensity:
\begin{equation}
  S_*(\Phi, \mathbf{v}) = \gamma\bigl(\Phi^2 + \|\mathbf{v}\|^2\bigr),
  \quad \gamma > 0.
  \label{eq:entropy_floor}
\end{equation}

This expression defines the minimal entropy required to encode the local
constraint complexity. Regions with large scalar density or strong flow require
greater entropy to remain admissible. This establishes a scaling relation
between structural intensity and informational overhead.

\begin{proposition}
  The entropy floor \eqref{eq:entropy_floor} is the unique quadratic function
  $S_*(\Phi,\mathbf{v}) = \gamma_1\Phi^2 + \gamma_2\|\mathbf{v}\|^2$ that is
  non-negative, vanishes at $(\Phi,\mathbf{v}) = (0,0)$, and is isotropic in
  $\mathbf{v}$. Setting $\gamma_1 = \gamma_2 = \gamma$ is the
  minimal-parameter choice consistent with equal weighting of scalar and vector
  burden.
\end{proposition}

\begin{proof}
  Non-negativity and vanishing at the origin require both coefficients to be
  non-negative. Isotropy in $\mathbf{v}$ requires the coefficient of
  $\|\mathbf{v}\|^2$ to be scalar. Equal weighting collapses two free
  parameters to one.
\end{proof}

\subsection{The Admissibility Potential}

Rather than enforcing \eqref{eq:admissibility} as a hard constraint, we
introduce a smooth penalty. Define
\begin{equation}
  V_{\mathrm{adm}}(S;\Phi,\mathbf{v})
  = \log\!\Bigl(1 + \exp\bigl(\alpha\,(S_*(\Phi,\mathbf{v}) - S)\bigr)\Bigr),
  \quad \alpha > 0.
  \label{eq:Vadm}
\end{equation}

This softplus barrier is smooth and differentiable everywhere. When $S$ is well
above the entropy floor, the potential is negligible. When $S$ falls below
$S_*$, the potential grows approximately linearly, imposing a strong penalty.
The parameter $\alpha$ controls the sharpness of the transition, while $\nu$
controls the strength of the penalty in the action.

\begin{proposition}[Properties of $V_{\mathrm{adm}}$]
  \label{prop:Vadm}
  The potential \eqref{eq:Vadm} satisfies: $V_{\mathrm{adm}} > 0$ for all
  finite arguments; $V_{\mathrm{adm}} \to 0$ as $S \to +\infty$;
  $V_{\mathrm{adm}} \sim \alpha(S_* - S)$ as $S \to -\infty$; and
  $\partial V_{\mathrm{adm}}/\partial S < 0$, so the penalty is strictly
  decreasing in entropy. Moreover $\partial V_{\mathrm{adm}}/\partial \alpha > 0$:
  larger $\alpha$ sharpens the transition at $S = S_*$.
\end{proposition}

\begin{proof}
  Direct computation from the softplus definition $\mathrm{sp}(u) =
  \log(1+e^u)$ with $u = \alpha(S_* - S)$.
  The derivative with respect to $S$ is
  $\partial_S V_{\mathrm{adm}} = -\alpha\,\sigma(\alpha(S_*-S))$,
  where $\sigma$ is the logistic function, which is strictly positive.
\end{proof}

This term gives entropy genuine regulatory power over the action. It does not
merely modulate coupling between $\Phi$ and $\mathbf{v}$; it determines whether
a configuration is admissible at all. Together with the entropy floor
\eqref{eq:entropy_floor}, it ensures that high field intensity demands high
entropy, driving the system toward a regime where structural complexity and
informational overhead are in balance.

\subsection{The Formal Action}

We write the full RSVP action as
\begin{align}
  \mathcal{A}[X]
  &= \int_\Omega
    \Bigl(
      \mathcal{L}_\Phi
      + \mathcal{L}_{\mathbf{v}}
      + \mathcal{L}_S
      + e^{-\lambda S}\,\mathcal{L}_{\Phi\mathbf{v}}
      + \mu\,(\nabla S \cdot \mathbf{J}_{\Phi,\mathbf{v}})
      + \nu\,V_{\mathrm{adm}}(S;\Phi,\mathbf{v})
    \Bigr)\,dx,
    \label{eq:full_action}
\end{align}
where $X = (\Phi,\mathbf{v},S)$. The action is composed of three classes of
terms. The first consists of standard single-field contributions governing local
magnitude and smoothness. The second consists of structural coupling between
$\Phi$ and $\mathbf{v}$, modulated by entropy. The third consists of
admissibility terms in which $S$ acts as a regulator.

\subsection{Structural Terms}

The one-field contributions are quadratic in the fields and their gradients:
\begin{align}
  \mathcal{L}_\Phi &= a_\Phi \Phi^2 + b_\Phi |\nabla\Phi|^2, \\
  \mathcal{L}_{\mathbf{v}} &= a_v \|\mathbf{v}\|^2 + b_v \|\nabla\mathbf{v}\|^2, \\
  \mathcal{L}_S &= a_S S^2 + b_S |\nabla S|^2.
\end{align}

These terms ensure that each field remains bounded and exhibits controlled
variation. The coupling between scalar and vector sectors is
\begin{equation}
  \mathcal{L}_{\Phi\mathbf{v}}
  = c_{\Phi v}\,\Phi\,(\nabla \cdot \mathbf{v})
    + c_{\Phi v}'\,|\nabla\Phi|^2\,\|\mathbf{v}\|^2.
  \label{eq:Lpv}
\end{equation}
The divergence term captures how scalar accumulation responds to compression and
expansion of flow; the gradient interaction term couples spatial variation in
$\Phi$ to the magnitude of $\mathbf{v}$.

\subsection{Entropy as Regulator: The Asymmetric Terms}

The defining feature of the RSVP action is the asymmetric role of entropy. The
factor $e^{-\lambda S}\,\mathcal{L}_{\Phi\mathbf{v}}$ modulates the
scalar--vector coupling. In regions of high entropy, coupling is suppressed; in
regions of low entropy, it is amplified. This is not a symmetric interaction.
It is a gating mechanism in which $S$ determines how strongly the structural
fields are allowed to interact.

The interaction current $\mathbf{J}_{\Phi,\mathbf{v}}$ is the variational
current of the coupling term:
\begin{equation}
  \mathbf{J}_{\Phi,\mathbf{v}}
  = \frac{\partial \mathcal{L}_{\Phi\mathbf{v}}}{\partial(\nabla\Phi)}
  + \frac{\partial \mathcal{L}_{\Phi\mathbf{v}}}{\partial(\nabla\mathbf{v})}.
  \label{eq:current}
\end{equation}
Evaluating from \eqref{eq:Lpv} with $c_{\Phi v}' = 0$, the current reduces to
$\mathbf{J}_{\Phi,\mathbf{v}} = \Phi\,\mathbf{v}$, the simplest
scalar-weighted flow vector. The term $\mu(\nabla S \cdot
\mathbf{J}_{\Phi,\mathbf{v}})$ penalizes configurations in which entropy
gradients are misaligned with the flow of scalar--vector coupling energy:
when entropy increases in the direction of structural transport, the
configuration is supported; when entropy gradients oppose the direction of
coupling, the configuration is penalized.

\subsection{Variational Equations of Motion}

Taking functional derivatives of \eqref{eq:full_action} yields the
Euler--Lagrange equations for each field. The entropy equation is the most
revealing:
\begin{equation}
  -b_S\,\Delta S + 2a_S\,S
  - \lambda\,e^{-\lambda S}\,\mathcal{L}_{\Phi\mathbf{v}}
  + \mu\,\nabla\cdot\mathbf{J}_{\Phi,\mathbf{v}}
  - \nu\,\alpha\,\sigma\!\bigl(\alpha(S_*-S)\bigr)
  = 0,
  \label{eq:EL_S}
\end{equation}
where $\sigma(u) = (1+e^{-u})^{-1}$ is the logistic function. The term
$-\lambda e^{-\lambda S}\mathcal{L}_{\Phi\mathbf{v}}$ drives $S$ upward where
coupling is strong, while $-\nu\alpha\sigma(\alpha(S_*-S))$ drives $S$ upward
when it falls below the entropy floor. Both forces push $S$ in the same
direction when the field is intense: the system self-regulates toward
admissibility. This is not imposed externally but emerges from the structure of
the action.

\subsection{Weight Functional}

The configuration weight is
\begin{equation}
  W[X] = e^{-\beta\,\mathcal{A}[X]}.
\end{equation}
Observables are computed as weighted averages over the configuration space.
However, direct evaluation of $W$ is exponentially intractable, as it requires
integration over all possible field configurations. The remainder of the essay
develops a method for constructing and compressing this weight through local
operations.

\begin{remark}[Double-exponential and linearization]
  When the gate is derived via $G_k \sim e^{-\beta\,\Delta\mathcal{A}_k}$, the
  entropy-gated coupling term $e^{-\lambda S}\mathcal{L}_{\Phi\mathbf{v}}$
  contributes $\exp(-\beta e^{-\lambda S}\Delta\mathcal{L}_{\Phi\mathbf{v}})$
  to the gate: a double exponential in $S$. This is analytically correct but
  numerically hazardous, as the dynamic range across local basis states can be
  extreme and degrades SVD truncation quality. The prototype scheme of
  Section~3 separates the admissibility and structural sectors via a Trotter
  split, avoiding this issue while preserving the conceptual asymmetry.
\end{remark}

% ============================================================
\section{The Prototype Decimation Scheme}
% ============================================================

\subsection{From Formal Action to Discrete Lattice}

To obtain a computationally tractable representation of the RSVP weight
functional, we discretize $\Omega$ as a one-dimensional chain of $N$ sites.
Each site $n$ carries a local state
\[
  x_n = (\Phi_n, v_n, S_n),
\]
drawn from a finite local basis of size $d = d_\Phi \times d_v \times d_S$.
The global weight $W(x_1,\dots,x_N) = e^{-\beta \mathcal{A}[X]}$ is not
assumed to factorize a priori. Instead, it is constructed through a sequence of
local operator applications derived from the discretized action. This
construction induces an approximate matrix product state (MPS) / tensor-train
representation~\cite{oseledets2011tt, schollwock2011dmrg}:
\begin{equation}
  W(x_1,\dots,x_N) \approx \sum_{\{\alpha\}}
  C_1^{\alpha_0\alpha_1}(x_1)\,C_2^{\alpha_1\alpha_2}(x_2)
  \cdots C_N^{\alpha_{N-1}\alpha_N}(x_N),
  \label{eq:MPS}
\end{equation}
where each core $C_n$ has shape $(\chi_{n-1}, d, \chi_n)$ and $\chi_0 =
\chi_N = 1$ (open boundary conditions). The bond dimension $\chi_n$ measures
retained correlations between neighboring sites.

The chain begins in a trivially factorized state, analogous to the
unentangled initial state of the Integral Decimation Method~\cite{grimm2026}.
All nontrivial dependence between sites is accumulated through the sequential
application of local operators. This representation emerges from the same
principle as the reference method: a global object is built incrementally
through local transformations, with compression applied at each step to
maintain tractability.

\subsection{Local Action Decomposition}

The discrete action decomposes as
\begin{equation}
  \mathcal{A}_{\mathrm{discrete}}[X]
  = \sum_n \mathcal{A}^{(1)}_n(x_n)
  + \sum_n \mathcal{A}^{(2)}_{n,n+1}(x_n, x_{n+1}),
\end{equation}
where the one-site term is
\begin{equation}
  \mathcal{A}^{(1)}_n(x_n)
  = a_\Phi\Phi_n^2 + a_v v_n^2 + a_S S_n^2,
  \label{eq:A1}
\end{equation}
and the two-site term is split into structural and admissibility sectors:
\begin{equation}
  \mathcal{A}^{(2)}_{n,n+1}(x_n,x_{n+1})
  = \mathcal{A}_{\Phi\mathbf{v}}(x_n,x_{n+1})
  + \mathcal{A}_{\mathrm{adm}}(x_n,x_{n+1}).
  \label{eq:A2}
\end{equation}

\subsection{Trotter Splitting of the Local Action}

The local contribution of the action on a pair of neighboring sites $(n,n+1)$
is denoted by $\Delta\mathcal{A}_{n,n+1}$. Rather than exponentiating the full
local action directly, we decompose it according to \eqref{eq:A2} and apply a
first-order Trotter splitting.

\begin{definition}[Trotter split]
  Given a two-site action of the form \eqref{eq:A2}, the
  \emph{first-order Trotter split} of the corresponding gate is
  \begin{equation}
    G_{n,n+1}
    = e^{-\beta\,\mathcal{A}^{(2)}_{n,n+1}}
    \approx G_{\mathrm{adm}}\,G_{\Phi\mathbf{v}}
    = e^{-\beta\,\mathcal{A}_{\mathrm{adm}}}\,
      e^{-\beta\,\mathcal{A}_{\Phi\mathbf{v}}}.
    \label{eq:trotter}
  \end{equation}
\end{definition}

\begin{proposition}[Trotter error and exact commutativity]
  \label{prop:trotter}
  The error in \eqref{eq:trotter} is bounded by
  $\frac{\beta^2}{2}\bigl\|[\mathcal{A}_{\Phi\mathbf{v}},
  \mathcal{A}_{\mathrm{adm}}]\bigr\|_F + O(\beta^3)$.
  In the prototype, $\mathcal{A}_{\mathrm{adm}}$ depends only on single-site
  $S$ values and $\mathcal{A}_{\Phi\mathbf{v}}$ depends only on $(\Phi,v)$
  pairs; on the local discrete basis these operators are simultaneously
  diagonal, so the commutator vanishes exactly and the Trotter split is exact.
\end{proposition}

\begin{proof}
  The Baker--Campbell--Hausdorff formula gives $e^{A+B} = e^A e^B
  e^{-[A,B]/2}\cdots$, so the leading error is $\frac{1}{2}[A,B]$ in the
  exponent. Since both gates are diagonal on the local basis, they commute on
  every basis element and the commutator is zero.
\end{proof}

The Trotter splitting therefore introduces zero first-order error in the
prototype regime, making it an exact operator factorization rather than an
approximation. This also preserves the conceptual separation between structure
and admissibility at the level of the numerical scheme.

\subsection{The Admissibility Gate}

The admissibility sector acts independently on each site within the pair:
\[
  \mathcal{A}_{\mathrm{adm}}(x,y)
  = \nu\,V_{\mathrm{adm}}(x) + \nu\,V_{\mathrm{adm}}(y),
\]
where $V_{\mathrm{adm}}$ is evaluated using the full local triplet
$(\Phi,v,S)$. The corresponding gate is diagonal:
\[
  G_{\mathrm{adm}}(x',y';x,y)
  = \delta_{x'x}\,\delta_{y'y}\,
    \exp\!\bigl(-\beta\,\mathcal{A}_{\mathrm{adm}}(x,y)\bigr).
\]
Because this gate depends only on local variables and does not couple
neighboring sites directly, it is computationally inexpensive. Its role is
purely regulatory: it suppresses configurations that violate the entropy floor
and leaves admissible configurations largely unchanged.

\subsection{The Structural Gate}

The structural sector captures the interaction between scalar and vector fields
across neighboring sites:
\[
  \mathcal{A}_{\Phi\mathbf{v}}(x,y)
  = b_\Phi(\Phi_x - \Phi_y)^2
  + b_v(v_x - v_y)^2
  + c_{\Phi v}\,\Phi_x\,v_y.
\]
This term enforces coherence between neighboring values while allowing directed
interaction through the asymmetric cross-term $\Phi_x v_y$. The corresponding
gate is also diagonal:
\[
  G_{\Phi\mathbf{v}}(x',y';x,y)
  = \delta_{x'x}\,\delta_{y'y}\,
    \exp\!\bigl(-\beta\,\mathcal{A}_{\Phi\mathbf{v}}(x,y)\bigr).
\]
Although diagonal, this gate induces correlations between neighboring sites by
reweighting configurations according to their structural compatibility.

\subsection{The Decimation Step}

After applying a two-site gate to cores $C_n$ and $C_{n+1}$, the update
proceeds in four steps.

\textbf{Step 1: Contract.}
\begin{equation}
  \Theta^{\ell,i,j,r} = \sum_m C_n^{\ell i m}\,C_{n+1}^{m j r}.
\end{equation}

\textbf{Step 2: Apply gate.}
Since the combined gate \eqref{eq:G_combined} is diagonal,
\begin{equation}
  \tilde\Theta^{\ell,i,j,r}
  = G_{n,n+1}(i,j;i,j)\,\Theta^{\ell,i,j,r}.
  \label{eq:G_combined}
\end{equation}

\textbf{Step 3: Reshape and SVD.}
\begin{equation}
  \Theta_{\mathrm{mat}} = \tilde\Theta.\mathrm{reshape}(Ld, dR),
  \qquad
  \Theta_{\mathrm{mat}} = U\Sigma V^\dagger.
\end{equation}

\textbf{Step 4: Truncate and split.}
Retain $k = \min\!\bigl(\max(1,\,|\{i:\sigma_i>\varepsilon\}|),\,
\chi_{\max}\bigr)$ singular values:
\begin{equation}
  C_n^{\mathrm{new}} =
    \bigl(U_{:,1:k}\,\diag(\sqrt{\sigma_{1:k}})\bigr)
    .\mathrm{reshape}(L,d,k),
  \quad
  C_{n+1}^{\mathrm{new}} =
    \bigl(\diag(\sqrt{\sigma_{1:k}})\,V^\dagger_{1:k,:}\bigr)
    .\mathrm{reshape}(k,d,R).
\end{equation}

The update rule is:
\begin{equation}
  W_{k+1} = \mathfrak{D}_\varepsilon(G_k\,W_k),
  \label{eq:update_rule}
\end{equation}
where $\mathfrak{D}_\varepsilon$ denotes the truncation operator and $G_k$ the
gate at bond $k$. This is the core decimation equation of the scheme. It
eliminates components of the representation that contribute negligibly to the
global weight, enforcing an entropy-respecting compression of the system.

\subsection{Sweep Protocol}

The full update proceeds through alternating sweeps. A left-to-right pass
applies one-site and two-site gates sequentially across all sites and bonds,
followed by a right-to-left pass in reverse order. After each pass, the tensor
train is brought into left-canonical form by QR decomposition, ensuring that
contractions used in observable computation remain well-conditioned. The full
sweep map is:
\begin{equation}
  \mathcal{T} = \mathfrak{D}_\varepsilon \circ G_{\mathrm{sweep}},
  \label{eq:sweep_map}
\end{equation}
where $G_{\mathrm{sweep}}$ denotes the full left-to-right-to-left gate
sequence. This mirrors the time-evolving block decimation (TEBD)
algorithm~\cite{vidal2004tebd}, adapted here to operate on a reweighted
configuration space rather than a unitary quantum state.

% ============================================================
\section{Implementation and Observables}
% ============================================================

\subsection{Simulator Architecture}

The simulator implements the decimation scheme as a composition of modular
components that mirror the mathematical structure of the method. The
\texttt{lattice} module defines the discrete domain, including the enumeration
of local states and the adjacency structure of the chain. The \texttt{core}
module stores the tensor-train representation, providing operations for
reshaping, contracting, and truncating cores. The \texttt{operators} module
constructs one-site and two-site gates from the Trotter-split action,
maintaining the separation between structural and admissibility contributions.
The \texttt{sweeps} module executes alternating passes, applying gates and
invoking truncation after each local update. The \texttt{observables} module
computes marginal distributions and derived quantities from the current
tensor-train state.

This modular decomposition reflects the separation of concerns present in the
theoretical construction. The action determines the operators, the operators
drive the evolution, and the tensor-train representation mediates between local
updates and global structure.

\subsection{Local Marginals and Field Expectations}

Left environments $L^{(n)} \in \mathbb{R}^{\chi_n \times \chi_n}$ are
computed recursively:
\begin{align}
  L^{(0)} &= [1], &
  L^{(n)}_{bb'} &= \sum_{a,a',i} L^{(n-1)}_{aa'}\,
    C_n^{a i b}\,C_n^{a' i b'}.
  \label{eq:left_env}
\end{align}
Right environments $R^{(n)}$ are computed analogously from the right boundary.
The local marginal probability at site $n$ is then:
\begin{equation}
  p_n(i)
  = \frac{1}{Z}
    \sum_{a,a',b,b'}
    L^{(n-1)}_{aa'}\,C_n^{a i b}\,C_n^{a' i b'}\,R^{(n)}_{bb'},
  \label{eq:marginal}
\end{equation}
where $Z = \sum_i p_n(i)$ before normalization. These environments are most
naturally computed in a canonical gauge, where orthogonality conditions prevent
uncontrolled growth or decay of intermediate quantities.

Field expectations at site $n$ are computed as
\begin{equation}
  \langle \Phi_n \rangle = \sum_i p_n(i)\,\Phi^{(i)},
  \qquad
  \langle v_n \rangle = \sum_i p_n(i)\,v^{(i)},
  \qquad
  \langle S_n \rangle = \sum_i p_n(i)\,S^{(i)}.
  \label{eq:expectations}
\end{equation}

\subsection{Bond Dimension as Closure Diagnostic}

The bond dimension $\chi_n$ at bond $(n,n+1)$ is the number of singular values
retained after truncation at that bond. It measures the amount of correlation
between neighboring sites that cannot be compressed below threshold
$\varepsilon$.

\begin{definition}[Effective bond dimension and closure burden]
  \begin{equation}
    \chi_{\mathrm{eff}}
    = \frac{1}{N-1}\sum_{n=1}^{N-1}\chi_n,
    \qquad
    \chi_{\max} = \max_{1\le n\le N-1}\chi_n.
    \label{eq:chi_eff}
  \end{equation}
  The \emph{closure burden} at bond $n$ is defined as $\chi_n / \chi_{\max}$.
  A uniform burden indicates homogeneous long-range correlation; localized
  peaks indicate sites where the field resists compression.
\end{definition}

A small $\chi_{\mathrm{eff}}$ indicates that the global configuration can be
expressed with low inter-site correlation, corresponding to a state of low
torsion and high admissibility. A large $\chi_{\max}$ localized at a particular
bond indicates a region of persistent structural interaction that resists
decimation, often associated with sharp field gradients or competing
constraints.

From the perspective of the framework, stabilization of bond dimensions under
repeated sweeps indicates that the system has reached a regime in which
constraint closure has been achieved at the chosen resolution. Information flow
across bonds becomes bounded, and the representation ceases to grow in
complexity. This boundedness is the empirical signature of admissibility.

\subsection{Simulation Results}

Figures \ref{fig:chi_beta} and \ref{fig:chi_nu} show $\chi_{\mathrm{eff}}$
per sweep for varying $\beta$ and $\nu$ respectively, on a chain of $N=20$
sites with $d = 27$ local states, $\chi_{\max} = 24$, and truncation threshold
$\varepsilon = 10^{-5}$. Bond dimension grows initially as correlations are
established, then stabilizes within approximately ten sweeps across all tested
parameter values. Figures \ref{fig:bonds} and \ref{fig:evolution} show the
spatial profile of bond dimensions at the final sweep and the sweep-by-sweep
evolution.

\begin{figure}[h]
\centering
\includegraphics[width=0.75\textwidth]{fig_chi_eff_by_beta.pdf}
\caption{Effective bond dimension $\chi_{\mathrm{eff}}$ per sweep for
  $\nu = 0.5$ and varying $\beta$. Stabilization within ten sweeps
  supports the contraction conjecture.}
\label{fig:chi_beta}
\end{figure}

\begin{figure}[h]
\centering
\includegraphics[width=0.75\textwidth]{fig_chi_eff_by_nu.pdf}
\caption{Effective bond dimension $\chi_{\mathrm{eff}}$ per sweep for
  $\beta = 1.0$ and varying $\nu$. Stronger admissibility coupling
  does not prevent stabilization.}
\label{fig:chi_nu}
\end{figure}

\begin{figure}[h]
\centering
\includegraphics[width=0.75\textwidth]{fig_bond_profile.pdf}
\caption{Closure burden profile at the final sweep ($\beta=1.0$,
  $\nu=0.5$). Localized peaks mark regions of persistent structural
  correlation.}
\label{fig:bonds}
\end{figure}

\begin{figure}[h]
\centering
\includegraphics[width=0.75\textwidth]{fig_bond_evolution.pdf}
\caption{Bond dimension evolution per sweep ($\beta=1.0$, $\nu=0.5$).
  Rows are sweeps; columns are bonds. Stabilization is visible within
  the first ten sweeps.}
\label{fig:evolution}
\end{figure}

% ============================================================
\section{The Convergence Conjecture}
% ============================================================

\subsection{The Truncated Sweep Map}

The full update procedure can be expressed as a map on the space of
tensor-train states. Let $\mathcal{M}_{\chi,N}$ denote the space of $N$-site
tensor trains with bond dimension at most $\chi$ in left-canonical form,
equipped with the Frobenius norm $\|W\|_F^2 = \sum_{x_1,\dots,x_N}
W(x_1,\dots,x_N)^2$.

Let $G$ denote the operator applying all gates in a complete sweep and
$\Pi_\varepsilon$ the nonlinear projection truncating the tensor train at
threshold $\varepsilon$. The truncated sweep map is
\begin{equation}
  \mathcal{T} = \Pi_\varepsilon \circ G.
  \label{eq:T}
\end{equation}

\subsection{Component Properties}

\begin{lemma}[Nonexpansiveness of $\Pi_\varepsilon$]
  \label{lem:nonexpansive}
  For any two tensor trains $W, W' \in \mathcal{M}_{\chi,N}$ in
  left-canonical form,
  \begin{equation}
    \|\Pi_\varepsilon(W) - \Pi_\varepsilon(W')\|_F \le \|W - W'\|_F.
  \end{equation}
\end{lemma}

\begin{proof}[Proof sketch]
  SVD truncation is the best rank-$k$ approximation in the Frobenius norm
  (Eckart-Young theorem). The map $W \mapsto \Pi_\varepsilon(W)$ is therefore
  a projection onto a closed set in the space of $(Ld)\times(dR)$ matrices.
  Projections onto closed convex sets are nonexpansive in Hilbert space norms.
  Applying this bond-by-bond preserves nonexpansiveness over the full chain.
\end{proof}

\begin{lemma}[Contractiveness of $G$ near identity]
  \label{lem:contractive}
  For sufficiently small $\beta$, the gate map $G$ satisfies
  $\|G(W) - G(W')\|_F \le (1 - c\beta)\|W - W'\|_F$
  for some $c > 0$ depending on the action coefficients.
\end{lemma}

\begin{proof}[Proof sketch]
  Each gate $G_k(i,j) = e^{-\beta\mathcal{A}(x^{(i)},x^{(j)})} \le 1$
  since $\mathcal{A} \ge 0$. The Frobenius norm therefore cannot increase under
  gate application. For strictly positive $\mathcal{A}$, at least some gate
  values are strictly less than 1, giving a strict contraction on the support.
\end{proof}

\subsection{Conjecture Statement}

\begin{conjecture}[Contraction and fixed point]
  \label{conj:main}
  Under the conditions of Lemmas~\ref{lem:nonexpansive} and
  \ref{lem:contractive}, the truncated sweep map $\mathcal{T} = \Pi_\varepsilon
  \circ G$ is contractive on $\mathcal{M}_{\chi_{\max},N}$. It has a unique
  fixed point $W^*$ satisfying
  \begin{equation}
    W^* = \mathcal{T}(W^*) = \Pi_\varepsilon(G(W^*)),
    \label{eq:fixed_point}
  \end{equation}
  which is the admissible compressed representation of the RSVP weight
  functional at parameters $(\beta, \varepsilon, \chi_{\max})$.
\end{conjecture}

\begin{remark}
  The main obstacle to a full proof is that $\Pi_\varepsilon$ is nonlinear,
  so standard Banach fixed-point arguments do not apply directly. A complete
  proof would require either showing $\Pi_\varepsilon$ is strictly contractive,
  or working in a metric space where the composition can be shown contractive
  directly. The empirical evidence supports the conjecture across all tested
  parameter regimes.
\end{remark}

The conjecture rests on two components. First, $\Pi_\varepsilon$ must be
nonexpansive in the chosen norm, which is expected when truncation error is
dominated by discarded singular values. Second, $G$ must be contractive, which
occurs when $\beta$ is small enough that the gates are close to identity and do
not amplify deviations.

\subsection{Empirical Support}

Across parameter regimes with $\beta \in [0.5, 2.0]$, $\varepsilon = 10^{-5}$,
$\chi_{\max} \in \{24\}$, and $\nu \in [0.2, 0.8]$, stabilization of
$\chi_{\mathrm{eff}}$ is observed within approximately ten sweeps (Figures
\ref{fig:chi_beta} and \ref{fig:chi_nu}). The stabilized value
$\chi_{\mathrm{eff}}^*$ decreases monotonically as $\varepsilon$ increases
or $\beta$ decreases, consistent with the expectation that contraction
strengthens as gate effects weaken and truncation becomes more aggressive. These
trends align with the predictions of Conjecture~\ref{conj:main}.

\subsection{Relation to Renormalization}

The truncated sweep map $\mathcal{T}$ bears a structural resemblance to
renormalization group transformations in statistical
mechanics~\cite{wilson1975rg, white1992dmrg}. In both cases, local
interactions are combined with a reduction of degrees of freedom to produce an
effective description of the system.

However, there is a crucial difference. Traditional renormalization proceeds
by integrating out short-range degrees of freedom to obtain a coarse-grained
theory at larger scales~\cite{wilson1975rg}. The decimation scheme presented
here preserves the full configuration space while compressing its
representation. No degrees of freedom are eliminated in principle; instead,
their contribution is encoded in a reduced set of correlations controlled by
the bond dimension.

The fixed point $W^*$ of $\mathcal{T}$ is therefore not an RG fixed point
--- a scale-invariant theory --- but a \emph{compression fixed point}: the most
compact tensor-train representation of the RSVP weight that survives repeated
entropy-respecting decimation at threshold $\varepsilon$. This is a new kind of
fixed point whose physical interpretation is precisely the notion of admissible
closure developed throughout this monograph.

% ============================================================
\section{Categorical Formulation of Constraint Closure}
% ============================================================

\subsection{The Category of Configurations}

We formalize the reconstruction process categorically. Let $\mathcal{C}$ be a
category whose objects are RSVP configurations represented as tensor-train
states and whose morphisms are admissible transformations induced by local
operators.

An object $W \in \mathrm{Ob}(\mathcal{C})$ is a tensor-train representation of
a weight functional. A morphism $f : W \to W'$ corresponds to the application
of a finite sequence of local gates followed by truncation. Composition of
morphisms corresponds to successive sweeps of the decimation process.

This category is not arbitrary. It is endowed with additional structure that
reflects the tensor-product decomposition of the system. In particular,
$\mathcal{C}$ carries a symmetric monoidal structure $(\mathcal{C}, \otimes, I)$
in which the tensor product corresponds to concatenation of independent
subsystems and $I$ is the trivial single-site system.

\subsection{The Admissibility Subcategory}

Define a full subcategory $\mathcal{C}_{\le \varepsilon} \subset \mathcal{C}$
consisting of tensor-train states whose bond dimensions are bounded by
$\chi_{\max}$ and whose truncation error is below $\varepsilon$ at every bond.
These objects represent admissible compressed configurations. The truncation
operator $\Pi_\varepsilon$ acts as a reflector:
\begin{equation}
  \Pi_\varepsilon : \mathcal{C} \to \mathcal{C}_{\le \varepsilon},
\end{equation}
mapping each $W \in \mathcal{C}$ to the closest admissible object in
$\mathcal{C}_{\le \varepsilon}$ with respect to the Frobenius norm.

\begin{proposition}
  $\Pi_\varepsilon$ is idempotent: $\Pi_\varepsilon \circ \Pi_\varepsilon =
  \Pi_\varepsilon$.
\end{proposition}

\begin{proof}
  Once a tensor-train state satisfies the truncation constraints, further
  application of $\Pi_\varepsilon$ does not alter it, since no singular values
  below threshold remain.
\end{proof}

\subsection{Constraint Closure as a Fixed Point}

Let $G : \mathcal{C} \to \mathcal{C}$ denote the endofunctor induced by a full
sweep of gate applications. The truncated sweep map is the composite endofunctor
$\mathcal{T} = \Pi_\varepsilon \circ G$, and constraint closure corresponds to
a fixed point:
\begin{equation}
  W^* \simeq \mathcal{T}(W^*).
\end{equation}

In categorical terms, $W^*$ is an algebra over the endofunctor $\mathcal{T}$:
a configuration invariant under the process of local generation followed by
admissibility filtering. This invariance expresses global consistency. The
configuration contains no components that would be removed by further
decimation, and therefore no components that fail to belong to a consistent
global section of the field.

\subsection{Relation to Sheaf-Theoretic Gluing}

The categorical formulation admits an interpretation in terms of sheaf theory.
Each local tensor core defines a section over a site, and the bond indices
encode compatibility conditions between neighboring sections. Constraint closure
corresponds to the existence of a global section obtained by gluing local data.
Failures of closure correspond to obstructions where local compatibility
conditions cannot be satisfied globally.

The truncation step $\Pi_\varepsilon$ can therefore be interpreted as removing
cohomological defects: components of the representation that prevent the
existence of a consistent global section. Decimation is a process of enforcing
sheaf consistency by eliminating non-gluable contributions. This connects to
the Baez--Stay Rosetta Stone~\cite{baezstay}, in which logical, physical, and
computational structures share a common categorical skeleton; the RSVP
decimation scheme realizes that skeleton concretely in a field-theoretic
setting.

% ============================================================
\section{Information-Theoretic Interpretation}
% ============================================================

\subsection{Entropy as Computational Resource}

The entropy field $S$ acquires a direct interpretation in information theory.
It measures the information required to encode a local configuration of
$(\Phi,\mathbf{v})$, and the entropy floor $S_*$ represents the minimal
information necessary to represent the local constraint structure.

The admissibility condition $S \ge S_*$ expresses a fundamental principle: no
structure can exist without sufficient information to encode it. This aligns
with Landauer's principle~\cite{landauer1961}, which relates information
processing to physical entropy, and with the broader statistical-mechanical
framework in which entropy constrains accessible states~\cite{jaynes1957}.
Bennett's analysis of reversible computation~\cite{bennett1982} and Lloyd's
bound on computational capacity~\cite{lloyd2000} further ground this
correspondence: the entropy floor $S_*(\Phi,\mathbf{v}) = \gamma(\Phi^2 +
\|\mathbf{v}\|^2)$ can be interpreted as a local Lloyd bound, setting the
minimum informational overhead required for a region of field intensity to
sustain a dynamical structure.

\subsection{Decimation as Information Compression}

The truncation operator $\Pi_\varepsilon$ performs lossy compression of the
tensor-train representation. Singular values below threshold $\varepsilon$
correspond to correlations that contribute negligibly to the global weight and
can be discarded without significantly altering the representation. This is
analogous to rate-distortion theory~\cite{coverthomas}, in which a signal is
compressed subject to a constraint on reconstruction error. Here, the
distortion measure is the Frobenius norm, and the rate is determined by the
bond dimension.

\begin{proposition}
  The bond dimension $\chi_n$ at bond $(n,n+1)$ provides an upper bound on the
  mutual information between the two halves of the system separated by that
  bond.
\end{proposition}

\begin{proof}
  In tensor-network representations, the entanglement entropy across a bond is
  bounded by $\log \chi_n$~\cite{schollwock2011dmrg}. Since mutual information
  is bounded by entanglement entropy, the result follows.
\end{proof}

Decimation therefore enforces an information constraint. Only correlations
representable within the available information budget are retained, reinforcing
the interpretation of constraint closure as an information-limited
reconstruction process.

\subsection{Connection to Predictive Models}

Predictive models optimize local likelihoods without enforcing global
constraints. From the present perspective, this corresponds to constructing
local marginals without ensuring they can be assembled into a consistent joint
distribution. The failure mode is well-known: locally coherent outputs that are
globally incoherent, a pathology not correctable by improving local accuracy
alone~\cite{vaswani2017, lecun2015}.

The tensor-train framework, by contrast, constructs a global object directly,
with local marginals derived from it. Decimation ensures that this global
object remains within an admissible complexity class. Robust intelligence
therefore requires mechanisms for maintaining global consistency under
information constraints, rather than optimizing local predictions in isolation.

% ============================================================
\section{Dynamical Systems Perspective on Constraint Closure}
% ============================================================

\subsection{State Space and Evolution Operator}

We reinterpret the truncated sweep map $\mathcal{T} = \Pi_\varepsilon \circ G$
as a discrete-time dynamical system acting on the state space
$\mathcal{M}_{\chi_{\max},N}$. Each iteration
\begin{equation}
  W_{k+1} = \mathcal{T}(W_k)
\end{equation}
defines a trajectory in this space. The evolution is composed of a smooth
transformation $G$ followed by a nonlinear projection $\Pi_\varepsilon$, placing
the system in the class of \emph{projected dynamical systems}, where evolution
is constrained to remain within an admissible subset.

\subsection{Invariant Sets and Attractors}

An invariant set $\mathcal{A} \subset \mathcal{M}_{\chi_{\max},N}$ satisfies
$\mathcal{T}(\mathcal{A}) \subseteq \mathcal{A}$. The fixed point $W^*$ of
Conjecture~\ref{conj:main} is a trivial invariant set. Empirically, no
non-trivial periodic or quasi-periodic behavior is observed in the prototype
regime; trajectories converge monotonically toward a stable configuration,
suggesting a global attractor consisting of a single fixed point.

\begin{proposition}
  If $\mathcal{T}$ is contractive, it admits a unique global attractor
  consisting of the single fixed point $W^*$~\cite{maclane1998}.
\end{proposition}

\begin{proof}
  Standard contraction mapping theory implies that all trajectories converge to
  the unique fixed point, and no other invariant sets can exist.
\end{proof}

The absence of oscillatory behavior is consistent with the interpretation of
decimation as an entropy-increasing process at the level of representation.
Each iteration removes incoherent components, reducing the effective
dimensionality of the state, and this monotonic reduction prevents the
recurrence of previous states.

\subsection{Lyapunov Functional}

Define
\begin{equation}
  \mathcal{L}(W) = \|W - \mathcal{T}(W)\|_F^2.
\end{equation}
This measures deviation from invariance under the update map and vanishes at the
fixed point.

\begin{proposition}
  Under the contraction assumption, $\mathcal{L}(W_k)$ is nonincreasing along
  trajectories.
\end{proposition}

\begin{proof}
  Contractiveness implies successive iterates approach $W^*$, reducing
  $\|W - \mathcal{T}(W)\|_F$ at each step.
\end{proof}

This serves as a Lyapunov function for the system, certifying stability of
the fixed point. In practice, $\mathcal{L}(W_k)$ can be estimated from the
change in $\chi_{\mathrm{eff}}$ between successive sweeps, providing a
convergence diagnostic computable without explicit knowledge of $W^*$.

\subsection{Entropy Production and Irreversibility}

The decimation step introduces irreversibility. Information discarded through
truncation cannot be recovered by subsequent iterations, distinguishing this
system from reversible Hamiltonian dynamics. This connects directly to
Landauer's principle~\cite{landauer1961}: the discarding of information has a
minimum thermodynamic cost of $k_B T \ln 2$ per bit, and the total truncation
error accumulated across sweeps represents a lower bound on the thermodynamic
cost of the computation.

The entropy field $S$ plays a dual role: it regulates admissibility in the
action, and the decimation process itself produces entropy at the level of
representation as information about discarded modes is irretrievably lost.
Both mechanisms act to eliminate configurations that cannot be sustained within
the available informational budget, creating a tight correspondence between
physical admissibility and computational irreversibility.

% ============================================================
\section{Geometric Interpretation}
% ============================================================

\subsection{Configuration Manifold}

The space of RSVP configurations can be viewed as a manifold $\mathcal{X}$
parameterized by $(\Phi,\mathbf{v},S)$. The action functional $\mathcal{A}[X]$
defines a geometry on this manifold, with lower-action regions corresponding to
more probable configurations under $W[X] = e^{-\beta\mathcal{A}[X]}$.

The admissibility condition defines a submanifold
\begin{equation}
  \mathcal{X}_{\mathrm{adm}}
  = \bigl\{ X \in \mathcal{X} \mid S(x) \ge S_*(\Phi(x),\mathbf{v}(x))
    \ \text{for a.e.}\ x \in \Omega \bigr\}.
\end{equation}
The soft potential $V_{\mathrm{adm}}$ smooths the boundary of this submanifold
into a transition region, ensuring that the action remains differentiable and
the Euler--Lagrange equations \eqref{eq:EL_S} are well-posed everywhere.

\subsection{Projection onto the Admissible Manifold}

The truncation operator $\Pi_\varepsilon$ acts geometrically as a projection
onto a low-dimensional manifold of admissible configurations within the space
of tensor-train representations. Each truncation step selects the closest point
in Frobenius norm within this manifold, removing components that lie outside
it. Repeated application drives the system toward the intersection of the image
of $G$ and the admissible manifold.

\subsection{Geodesic Interpretation of Sweeps}

Gate application $G$ can be interpreted as moving along a trajectory in
configuration space determined by local gradients of the action. In the
continuous limit, this corresponds to gradient flow:
\begin{equation}
  \frac{dX}{dt} = -\nabla \mathcal{A}[X].
\end{equation}
The discrete sweep process approximates this flow, while the truncation step
projects the trajectory back onto the admissible manifold. The overall dynamics
is therefore a \emph{projected gradient flow}:
\begin{equation}
  X_{k+1} = \Pi_{\mathrm{adm}}\!\bigl(X_k - \eta\,\nabla \mathcal{A}[X_k]\bigr),
  \label{eq:proj_grad}
\end{equation}
connecting the tensor-network algorithm to variational optimization methods
with admissibility enforced at each step rather than only at convergence.

\subsection{Curvature and Complexity}

Regions of high bond dimension correspond to regions of high curvature in the
configuration manifold, where local constraints interact complexly and require
many degrees of freedom to represent accurately. Decimation flattens these
regions by removing negligible directions of curvature.

The stabilized configuration lies in a region of reduced curvature, where the
manifold is effectively lower-dimensional. Constraint closure is therefore the
process of finding a low-curvature submanifold approximating the full
configuration space while preserving admissible structure. The bond-dimension
profile of Figure~\ref{fig:bonds} is a direct measurement of this curvature
distribution along the lattice.

% ============================================================
\section{Extensions and Open Problems}
% ============================================================

\subsection{Higher-Dimensional Lattices}

The prototype scheme is formulated on a one-dimensional chain. Extension to
higher-dimensional lattices requires replacing tensor trains with projected
entangled pair states (PEPS) or tree tensor networks~\cite{orus2014tensor}.
The main challenge is computational tractability, as exact contraction costs
grow exponentially with dimension. Nevertheless, the conceptual framework
remains unchanged: local operators are applied to neighboring regions,
decimation controls correlation growth, and the admissibility potential
$V_{\mathrm{adm}}$ requires no modification.

\subsection{Continuous Limits and Field Reconstruction}

A key open problem is the recovery of continuous field behavior from the
discrete tensor-train representation. In the limit of large $N$ and refined
local bases, the discrete model should approximate solutions to the
Euler--Lagrange equations derived from the RSVP action. Establishing this
correspondence rigorously would connect the computational scheme to the
underlying field theory, analogous to the relationship between lattice field
theory and continuum quantum field theory.

\subsection{Adaptive Truncation and Dynamic Resolution}

The truncation threshold $\varepsilon$ and maximum bond dimension $\chi_{\max}$
are currently fixed parameters. An adaptive scheme varying these quantities
across the lattice would improve efficiency by allocating more resources to
regions of high closure burden. This aligns with the interpretation of bond
dimension as a curvature measure: dynamically adjusting computational effort
according to local constraint complexity.

\subsection{Toward a Proof of Convergence}

The convergence conjecture remains the central open theoretical problem.
Possible approaches include: constructing a weighted norm in which
$\Pi_\varepsilon$ becomes strictly contractive; bounding the Lipschitz constant
of $G$ in terms of action parameters $(\beta, a_\Phi, a_v, a_S, b_\Phi, b_v,
\nu)$; or reformulating in terms of monotone operator theory, which accommodates
nonlinear projections more naturally than classical Banach fixed-point
arguments. Each pathway would provide a rigorous foundation for the empirical
observations presented here.

\subsection{Cognitive and Physical Implications}

The framework suggests a reinterpretation of both cognition and physical
dynamics. The Friston free-energy principle~\cite{friston2010} frames
perception as minimization of a variational bound on surprise; constraint
closure provides a complementary picture in which consistency, rather than
surprise reduction, is the organizing criterion. Thought processes may be
understood as iterative construction of admissible configurations, with
inconsistencies eliminated through a process analogous to decimation.

In physics, the Jacobson derivation of Einstein's equations from thermodynamic
principles~\cite{jacobson1995} and Verlinde's entropic gravity
program~\cite{verlinde2011} both locate gravitational structure in entropic
constraints rather than metric dynamics. The RSVP framework is consistent with
this perspective and extends it: the entropy floor $S_*$ connects local field
intensity to informational overhead in a manner that could, in principle, be
extended to a cosmological setting in which spacetime structure emerges from
constraint closure rather than from a background metric.

% ============================================================
\section{Conclusion: Constraint Closure as a Unifying Principle}
% ============================================================

The development presented in this essay establishes a continuous chain from
formal field theory to concrete computational implementation and finally to an
empirically grounded conjecture about convergence. At each stage, the same
structural principle appears: configurations are generated through local
interactions, but only those that satisfy global admissibility constraints are
retained. The mechanism by which this selection occurs is entropy-respecting
decimation.

Within the RSVP framework, this principle is encoded directly in the action.
The scalar and vector fields generate structure through local interactions,
while the entropy field regulates whether that structure can be sustained. The
admissibility potential enforces a relationship between local intensity and
informational overhead, ensuring that configurations lacking sufficient entropy
are suppressed. This introduces a hierarchy in which structure is subordinate
to admissibility, rather than the two being treated as coequal components of
the dynamics. As the entropy equation of motion \eqref{eq:EL_S} reveals, this
hierarchy is not imposed externally but emerges from the variational structure
of the action itself.

When the action is discretized and translated into a computational scheme, this
same hierarchy persists. The Trotter-split decomposition separates structural
and admissibility effects into distinct operators, with Proposition
\ref{prop:trotter} establishing that this split is exact on the local basis.
The sequential application of these operators constructs a global
representation through local updates. The decimation step then enforces
admissibility by removing components that cannot be integrated into a coherent
structure. The tensor-train representation serves as the medium through which
this process is realized, allowing global structure to emerge from local
operations while remaining computationally tractable.

The convergence conjecture formalizes the observed behavior of this system.
Stabilization of bond dimension under repeated sweeps indicates that the
representation has reached a regime in which further application of $\mathcal{T}$
does not increase complexity. Lemmas \ref{lem:nonexpansive} and
\ref{lem:contractive} establish the two components required for contraction,
and Conjecture \ref{conj:main} states the expected fixed-point result.
The categorical formulation of Section~6 shows that the fixed point is an
algebra over the endofunctor $\mathcal{T}$ and that decimation enforces sheaf
consistency by removing non-gluable contributions. The information-theoretic
analysis of Section~7 connects the entropy floor to Landauer's principle and
frames bond dimension as a mutual-information bound. The dynamical-systems
perspective of Section~8 supplies a Lyapunov functional and identifies the
irreversibility of truncation with thermodynamic entropy production. The
geometric interpretation of Section~9 recasts sweeps as projected gradient
flow and bond-dimension peaks as regions of high manifold curvature.
Although a formal proof of Conjecture~\ref{conj:main} remains open, the
convergence of these four perspectives provides strong evidence that
constraint closure is not merely a conceptual idea but a measurable and
geometrically grounded property of the dynamics.

The implications of this framework extend beyond the specific numerical method.
In computation, it suggests that the limitations of predictive models arise
from their failure to enforce global consistency, and that more robust systems
will require mechanisms for constraint closure. In cognition, it suggests that
understanding may be better modeled as the process of assembling admissible
configurations rather than predicting isolated outcomes. In physics, it
provides a perspective in which structure arises from the selective retention
of configurations that satisfy both local interactions and global admissibility
conditions.

The unifying claim is therefore that stable structure --- whether physical,
computational, or cognitive --- is the result of a two-stage process. First,
local operators generate candidate configurations through the accumulation of
constraints. Second, a decimation mechanism filters these configurations
according to an admissibility criterion, removing those that cannot be
integrated into a coherent whole. The combination of these processes yields a
system that is both generative and selective, capable of producing complex
structure while maintaining global consistency.

Future work will refine each component of this framework. On the formal side,
the variational structure of the RSVP action can be extended to incorporate
additional fields and symmetries, and the role of the entropy field can be
explored in higher-dimensional settings. On the computational side, the
tensor-train implementation can be generalized to higher-dimensional lattices
and improved through adaptive truncation strategies. On the theoretical side,
the convergence conjecture can be investigated more rigorously, with the goal
of establishing conditions under which $\mathcal{T}$ is provably contractive.
The broader objective is a unified theory in which computation, cognition, and
physical dynamics are understood as manifestations of the same underlying
process of constraint closure.

\newpage

\begin{thebibliography}{99}

\bibitem{grimm2026}
R. T. Grimm, A. J. Staat, and J. D. Eaves,
\textit{The Integral Decimation Method for Quantum Dynamics and Statistical Mechanics},
Quantum, 10, 2064 (2026).
DOI: \href{https://doi.org/10.22331/q-2026-04-13-2064}{10.22331/q-2026-04-13-2064}.

\bibitem{oseledets2011tt}
I. V. Oseledets,
\textit{Tensor-Train Decomposition},
SIAM Journal on Scientific Computing, 33(5), 2295--2317 (2011).

\bibitem{orus2014tensor}
R. Or\'{u}s,
\textit{A Practical Introduction to Tensor Networks: Matrix Product States
and Projected Entangled Pair States},
Annals of Physics, 349, 117--158 (2014).

\bibitem{schollwock2011dmrg}
U. Schollw\"{o}ck,
\textit{The Density-Matrix Renormalization Group in the Age of Matrix
Product States},
Annals of Physics, 326(1), 96--192 (2011).

\bibitem{white1992dmrg}
S. R. White,
\textit{Density Matrix Formulation for Quantum Renormalization Groups},
Physical Review Letters, 69(19), 2863--2866 (1992).

\bibitem{vidal2004tebd}
G. Vidal,
\textit{Efficient Simulation of One-Dimensional Quantum Many-Body Systems},
Physical Review Letters, 93(4), 040502 (2004).

\bibitem{wilson1975rg}
K. G. Wilson,
\textit{The Renormalization Group: Critical Phenomena and the Kondo Problem},
Reviews of Modern Physics, 47(4), 773--840 (1975).

\bibitem{cichocki2016}
A. Cichocki et al.,
\textit{Tensor Networks for Dimensionality Reduction and Large-Scale
Optimization},
Foundations and Trends in Machine Learning, 9(4-5), 249--429 (2016).

\bibitem{hackbusch2012}
W. Hackbusch,
\textit{Tensor Spaces and Numerical Tensor Calculus},
Springer (2012).

\bibitem{landauer1961}
R. Landauer,
\textit{Irreversibility and Heat Generation in the Computing Process},
IBM Journal of Research and Development, 5(3), 183--191 (1961).

\bibitem{bennett1982}
C. H. Bennett,
\textit{The Thermodynamics of Computation---A Review},
International Journal of Theoretical Physics, 21(12), 905--940 (1982).

\bibitem{lloyd2000}
S. Lloyd,
\textit{Ultimate Physical Limits to Computation},
Nature, 406, 1047--1054 (2000).

\bibitem{jaynes1957}
E. T. Jaynes,
\textit{Information Theory and Statistical Mechanics},
Physical Review, 106(4), 620--630 (1957).

\bibitem{coverthomas}
T. M. Cover and J. A. Thomas,
\textit{Elements of Information Theory},
Wiley (2006).

\bibitem{jacobson1995}
T. Jacobson,
\textit{Thermodynamics of Spacetime: The Einstein Equation of State},
Physical Review Letters, 75(7), 1260--1263 (1995).

\bibitem{verlinde2011}
E. Verlinde,
\textit{On the Origin of Gravity and the Laws of Newton},
Journal of High Energy Physics, 2011(4), 29 (2011).

\bibitem{friston2010}
K. Friston,
\textit{The Free-Energy Principle: A Unified Brain Theory?},
Nature Reviews Neuroscience, 11(2), 127--138 (2010).

\bibitem{vaswani2017}
A. Vaswani et al.,
\textit{Attention Is All You Need},
NeurIPS (2017).

\bibitem{lecun2015}
Y. LeCun, Y. Bengio, and G. Hinton,
\textit{Deep Learning},
Nature, 521, 436--444 (2015).

\bibitem{goodfellow2014}
I. Goodfellow et al.,
\textit{Generative Adversarial Nets},
NeurIPS (2014).

\bibitem{maclane1998}
S. Mac Lane,
\textit{Categories for the Working Mathematician},
Springer (1998).

\bibitem{lurie2009}
J. Lurie,
\textit{Higher Topos Theory},
Princeton University Press (2009).

\bibitem{baezstay}
J. C. Baez and M. Stay,
\textit{Physics, Topology, Logic and Computation: A Rosetta Stone},
In New Structures for Physics, Springer (2010).

\bibitem{turing1936}
A. M. Turing,
\textit{On Computable Numbers, with an Application to the
Entscheidungsproblem},
Proceedings of the London Mathematical Society, 42(2), 230--265 (1936).

\bibitem{church1936}
A. Church,
\textit{An Unsolvable Problem of Elementary Number Theory},
American Journal of Mathematics, 58(2), 345--363 (1936).

\bibitem{shannon1948}
C. E. Shannon,
\textit{A Mathematical Theory of Communication},
Bell System Technical Journal, 27, 379--423 (1948).

\end{thebibliography}

\end{document}
